\documentclass[12pt]{article}
\usepackage[utf8]{inputenc}
\usepackage[T1]{fontenc}
\usepackage{hyperref}
\title{Thulung Language Resources for K-12 Curriculum Development}
\author{Compiled by Claude Sonnet 5 (claude-sonnet-5) for Andr\'as Kornai}
\date{}
\begin{document}
\maketitle
\section{Introduction}
Thulung (ISO 639-3 tdh; Glottocode thul1246) is a Sino-Tibetan (Kiranti) language of Nepal. The reference description gives about 21,000 speakers. Searches found documentation aimed at linguists, but no school materials.
\section{Primary Educational Resources}
No primers, textbooks or school programmes were found in search results. Thulung is not traditionally written and has no standard orthography.
\section{Linguistic and Reference Materials}
\begin{itemize}
\item Aspects of the Grammar of Thulung Rai: an Endangered Himalayan Language (Lahaussois, UC Berkeley dissertation, 2002): \url{https://lx.berkeley.edu/publications/aspects-grammar-thulung-rai-endangered-himalayan-language}
\item Language Description Heritage entry for Thulung: \url{https://ldh.clld.org/languages/tdh}
\item Huma-Num (CoCoon) catalogue entry, which lists a grammar of more than 300 pages: \url{https://cocoon.huma-num.fr/exist/crdo/lexvo.xql?code=tdh}
\item Wikipedia overview: \url{https://en.wikipedia.org/wiki/Thulung_language}
\end{itemize}
Allen's \emph{Sketch of Thulung Grammar} (1975) was named in the search summary, but no direct URL for it appeared in the results, so none is given here. No modern dictionary was found.
\section{Religious and Community Resources}
None found in search results.
\section{Digital and Audio Resources}
The Huma-Num catalogue entry above is the only archive-type listing found. Endangered Languages Project entry: \url{https://endangeredlanguages.com/lang/5101}
\section{Curriculum Development Context}
Little exists beyond academic grammars. A native-language curriculum would need an orthography decision, a lexicon and community-led materials. This would build on Lahaussois's description.
\end{document}